% \newglossaryentry{key}{
% name={},
% % type={},
% % text={},    % default = name
% % plural={},  
% % first={},   % default = text
% % firstplural={},
% % symbol={},
% % symbolplural={},
% % sort={},
% description={},
% % see={},
% }


\newglossaryentry{Asthma_bronchiale}{
name={Asthma bronchiale},
% type={},
% text={},    % default = name
% plural={},
% first={},   % default = text
% firstplural={},
% symbol={},
% symbolplural={},
% sort={},
description={Chronische Erkrankung mit mehr oder weniger
häufigen Episoden von Anfällen mit massiver Atemnot oder
Atembehinderung durch Verengung der Bronchien.},
% see={},
}

\newglossaryentry{Bewusstsein}{
name={Bewusstsein},
% type={},
% text={},    % default = name
% plural={},
% first={},   % default = text
% firstplural={},
% symbol={},
% symbolplural={},
% sort={},
description={Überbegriff u.a. für Wachheit,
Orientierung, Aufmerksamkeit, Auffassungsgabe, Denkverlauf und
Merkfähigkeit\cite{Pschy259}. \Z{Die "`wahre"' Definition
von Bewusstsein ist seit tausenden von Jahren ungeklärt. Auch an dieser Stelle
wird das Geheimnis nicht gelüftet werden, daher muss die obige Definition für
den Alltag ausreichen.\protect\footnote{Hinweisen zufolge könnte die Antwort
eventuell "`42"' sein.}}},
% see={},
}

\newglossaryentry{blutdruck}{
name={Blutdruck},
description={Durch die Pumpfunktion des Herzens entsteht ein
lebenswichtiger Druck im (arteriellen) Gefäßsystem, der
\E{Blutdruck}},
}

\newglossaryentry{Entzuendung}{
name={Entzündung},
% type={},
% text={},    % default = name
plural={Entzündungen},
% first={},   % default = text
% firstplural={},
% symbol={},
% symbolplural={},
% sort={},
description=
{
Reaktion des Körpers welche charakterisiert
ist durch  \Es{R\"otung}, \Es{Schwellung}, \E{\"Uberw\"armung},  \Es{Schmerz}
und \Es{Funktionsverlust}. \protect\Z{\protect\Lang{engl., lat.} inflammation.}
},
% see={},
}

\newglossaryentry{Herzinfarkt}{
name={Herzinfarkt},
description={Verschluss eines oder mehrerer Herzkranzgefäße mit
anschließendem Absterben des betroffenen Herzmuskelgewebes.},
% plural={Katastrophen}
}

\newglossaryentry{katastrophe}{
name={Katastrophe},
description={\protect\Lang{griech.}: schweres Unglück, Zusammenbruch,
Wendung zum Schlimmen \cite{BertelsmannVolkslexikon25,BrockhausDrei2Bd2}.},
plural={Katastrophen}
}

\newglossaryentry{Millimeter_Quecksilbersaeule}{
name={Millimeter Quecksilbersäule},
description={\Lang{Abkz.} mm\,Hg, auch: mmHg. Alte Maßeinheit für Druck, heute
Verwendung zur Angabe des Blutdruckes. \Z{Die Einheit ist definiert als der
hydrostatische Druck, der von 1\,mm Quecksilber ausgeübt wird.}},
symbol={mm\,Hg} }

\newglossaryentry{neurologie}{
name={Neurologie},
description={\Lang{allgem.} Lehre vom Nervensystem; \Lang{med.} Sonderfach
betreffend Erkrankungen des Nervensystems.}, }

\newglossaryentry{Notfallpatient}{
name={Notfallpatient},
% type={},
% text={},    % default = name
plural={Notfallpatienten},
% first={},   % default = text
% firstplural={},
% symbol={},
% symbolplural={},
% sort={},
description={Notfallpatienten gemäß  sind Patienten, bei denen im Rahmen einer
akuten Erkrankung, einer Vergiftung oder eines Traumas eine lebensbedrohliche
Störung einer vitalen Funktion eingetreten ist, einzutreten droht oder nicht
sicher auszuschließen ist.\protect\footnote{Abs. 1 Z 2 SanG}},
% see={},
}

\newglossaryentry{Schock}{
name={Schock},
description={Lebensbedrohliche Kreislaufstörung aufgrund von
Unterversorgung des Körpers mit Blut und Sauerstoff infolge des
Versagens des Kreislaufsystems.},
% plural={Katastrophen}
}